\documentclass[11pt]{article}
\usepackage{amsmath,amssymb}
\usepackage{geometry}
\usepackage{enumitem}
\geometry{margin=0.8in}
\setlength{\parindent}{0pt}

\begin{document}

\begin{center}
  {\Large\textbf{When Is a Rule a Function?}}\\[3pt]
  {\large Algebra II Word Problems and Abstract Reasoning}
\end{center}

Name: \rule{2.7in}{0.4pt}\hfill Date: \rule{1.5in}{0.4pt}

\section*{Instructions}
Suggested time: 45 minutes. Complete all 12 problems without a calculator.
You may use this reference page. For every function decision, identify the
input, the proposed output, and why the rule does or does not meet the
definition. Support a rejection with a specific input or counterexample.
Keep the stated domain and codomain unless a question asks you to change them.

\begingroup
\small
\section*{The Definition}
A function $f:D\to C$ assigns \textbf{exactly one element of $C$ to every
element of $D$}.
\begin{itemize}[leftmargin=*,itemsep=0.3em]
  \item \textbf{Domain $D$:} the full set of inputs the rule promises to accept.
  \item \textbf{Codomain $C$:} the declared set in which outputs must lie.
  \item \textbf{Range:} the outputs actually obtained,
  $\{f(x):x\in D\}\subseteq C$.
\end{itemize}

\section*{Three Requirements}
For each input in the stated domain:
\begin{enumerate}[leftmargin=*,itemsep=0.3em]
  \item An output must exist.
  \item The output must be unique.
  \item The output must belong to the stated codomain.
\end{enumerate}
A rule can fail for different reasons at different inputs. Changing its domain,
codomain, or assignment rule changes the proposed function.

\section*{What Is Allowed?}
\begin{itemize}[leftmargin=*,itemsep=0.3em]
  \item Different inputs may have the same output. The function need not use
  every element of its codomain.
  \item Inputs and outputs need not be numbers. They may be names, ordered
  pairs, sets, or even functions. An entire set can be one output object.
  \item A function is \textbf{one-to-one} if different inputs always have
  different outputs. This is an additional property, not part of being a function.
  \item If the same input object can be represented in different ways,
  its output must not depend on which representation is used.
\end{itemize}

\section*{Relations and Composition}
A relation from $D$ to $C$ is a collection of ordered pairs $(x,y)$ with
$x\in D$ and $y\in C$. It represents a function on $D$ only if each element
of $D$ appears as a first coordinate with exactly one second coordinate.

\medskip
For $(g\circ f)(x)=g(f(x))$ to be defined on all of $D$, $f$ must be defined
on $D$ and every resulting $f(x)$ must be an allowed input to $g$.
\endgroup

\newpage
\begin{center}
  {\large\textbf{When Is a Rule a Function? --- Practice}}
\end{center}

\begin{enumerate}[leftmargin=*,itemsep=0pt]
  \item A school has student IDs $S=\{A,B,C,D\}$ and lockers
  $L=\{1,2,3,4,5\}$. Its complete assignment record is
  \[
    A\mapsto2,\qquad B\mapsto2,\qquad C\mapsto4,\qquad D\mapsto5.
  \]
  Does the record define a function from $S$ to $L$? State its range and
  whether it is one-to-one. A second rule takes any locker in $L$ and
  returns the ID of a student assigned to it. Is that a function from
  $L$ to $S$? Identify every type of failure that occurs.
  \vspace{\stretch{1}}

  \item A library has a fixed inventory $B$ of book IDs and a set $S$ of
  registered students. At noon, one student has no borrowed books, another
  has two different books, and every loan is recorded. Compare these rules:
  \begin{enumerate}[label=(\alph*),leftmargin=1.6em,itemsep=0.25em]
    \item Input a student in $S$; output a book ID in $B$ that the student
    has borrowed.
    \item Input a student in $S$; output the complete set of that student's
    borrowed book IDs. The codomain is the collection of all subsets of $B$,
    including the empty set.
  \end{enumerate}
  Which rules are functions? Explain what happens for both described students.
  \vspace{\stretch{1}}

  \item Two weather stations each report at 8 a.m. and 8 p.m. on one fixed
  date. North reports $6^\circ$ and $11^\circ$; South reports $4^\circ$ and
  $4^\circ$. A rule takes a station name and returns a reported temperature
  from that date. Does this give a function from station names to real
  temperatures? Would using the ordered pair (station, reporting time) as
  the input repair the rule? State the full new domain and explain.
  \vspace*{\stretch{1}}
\end{enumerate}

\newpage
\begin{center}
  {\large\textbf{When Is a Rule a Function? --- Continued}}
\end{center}

\begin{enumerate}[leftmargin=*,itemsep=0pt,resume]
  \item A courier accepts parcel masses $m\in(0,5]$ kilograms. Its fee chart
  says ``\$5 when $0<m\le2$'' and ``\$8 when $2\le m\le5$.'' Does the chart
  define a function from accepted masses to dollar fees? Repair it by
  changing exactly one endpoint inequality while keeping the full domain.
  If both listed fees had instead been \$5, would the original overlap
  prevent the chart from defining a function? Explain.
  \vspace{\stretch{1}}

  \item A design program accepts a positive area $A$ and returns the
  perimeter of a rectangle having that area, with no further restrictions.
  Does this define a function from positive real areas to positive real
  perimeters? Give two specific rectangles to support your answer. If the
  program is restricted to squares, does the rule become a function?
  Give its formula and its domain and codomain.
  \vspace{\stretch{1}}

  \item An integer machine accepts any integer $n$ and must output an
  integer $m$ satisfying $2m=n$. Does this define a function
  $\mathbb Z\to\mathbb Z$? Explain what happens for even and odd inputs.
  Give two separate repairs: one that changes only the domain and keeps
  infinitely many inputs, and one that changes only the codomain while
  keeping every integer input. State both repaired functions precisely.
  \vspace*{\stretch{1}}
\end{enumerate}

\newpage
\begin{center}
  {\large\textbf{When Is a Rule a Function? --- Continued}}
\end{center}

\begin{enumerate}[leftmargin=*,itemsep=0pt,resume]
  \item A machine accepts a real number $x$ and returns a real number $y$
  satisfying $y^2=x$. Determine what happens when $x<0$, $x=0$, and $x>0$.
  Does restricting inputs to $x\ge0$ alone make the rule a function into
  $\mathbb R$? Give a domain and an output restriction that produce a
  function with infinitely many inputs, and write its rule.
  \vspace{\stretch{1}}

  \item A program takes a rational number $r$. It writes $r=p/q$, where
  $p$ is an integer and $q$ is a positive integer, and outputs $p+q$.
  Fractions need not be in lowest terms. Does this define a function
  $\mathbb Q\to\mathbb Z$? Use the input $r=1/2$ to justify your answer.
  Would the assignment $p+q$ define a function if the input were instead
  the ordered pair $(p,q)$? State that new domain and explain the difference.
  \vspace{\stretch{1}}

  \item A program accepts any nonempty subset $S$ of $\{1,2,3\}$. A set
  has no intrinsic order: $\{1,2\}=\{2,1\}$. For each proposed output below,
  decide whether the rule defines a function into $\{0,1,2,3\}$.
  For each valid function, state its range.
  \begin{enumerate}[label=(\alph*),leftmargin=1.6em,itemsep=0.25em]
    \item The number of elements of $S$.
    \item The first number written when someone lists the elements of $S$.
    \item The smallest element of $S$.
  \end{enumerate}
  \vspace*{\stretch{1}}
\end{enumerate}

\newpage
\begin{center}
  {\large\textbf{When Is a Rule a Function? --- Continued}}
\end{center}

\begin{enumerate}[leftmargin=*,itemsep=0pt,resume]
  \item A program's input is an entire real quadratic polynomial
  $p(x)=Ax^2+Bx+C$, with $A\ne0$. One version returns the real number $p(2)$.
  Does it define a function from quadratic polynomials to $\mathbb R$?
  Is it one-to-one? Exhibit two different inputs to support your answer.
  Another version returns a real zero of the input polynomial. Does that
  version define a function on the same domain? Give separate examples
  showing every type of failure that occurs.
  \vspace{\stretch{1}}

  \item Three students have advisor assignments
  $S_1\mapsto A_1$, $S_2\mapsto A_1$, and $S_3\mapsto A_2$.
  The advisor directory contains $A_1$, $A_2$, and $A_3$.
  Advisors $A_1$ and $A_2$ both have room 101; $A_3$ has no assigned room.
  Consider the student-to-advisor rule and the advisor-to-room rule, using
  all listed students and all listed advisors as their respective domains.
  Is each rule a function? Does ``look up a student's advisor, then that
  advisor's room'' nevertheless define a function on all three students?
  Explain, and state a restriction of the advisor domain that makes the
  second rule a function and permits the composition.
  \vspace{\stretch{1}}

  \item Invent a real-world assignment with exactly four distinct input
  objects, a codomain of exactly four distinct output labels, and a range
  of exactly three labels. List all input--output pairs and explain why
  your rule is a function but is not one-to-one. Then alter the relation,
  keeping the same stated domain and codomain, so that it is no longer a
  function. Identify the precise requirement that your altered relation fails.
  \vspace*{\stretch{1}}
\end{enumerate}

\end{document}
